% TOP_PROBLEM: {"id": 30001233, "title": "SHGH Conjecture for Linear Systems on Blown-Up Planes", "category_id": 5, "category": {"id": 5, "name": "algebraic_geometry", "display_name": "Algebraic Geometry", "description": "Geometric objects defined by polynomial equations.", "slug": "algebraic-geometry", "order_index": 5, "created_at": "2026-07-31T15:26:25.671Z"}, "status": "open"}
% FIRST_POSTED: 2026-09-27
% !TeX program = lualatex
% ATTEMPT_STATUS: unresolved
\documentclass[11pt]{article}
\usepackage[margin=25mm]{geometry}
\usepackage{fontspec}
\setmainfont{Latin Modern Roman}
\usepackage{amsmath,amssymb,mathtools}
\usepackage{unicode-math}
\setmathfont{Latin Modern Math}
\usepackage{xurl,hyperref}
\hypersetup{colorlinks=true,urlcolor=blue}
\setcounter{secnumdepth}{0}
\setlength{\parindent}{0pt}
\setlength{\parskip}{5pt}
\setlength{\emergencystretch}{3em}
\begin{document}
\section*{\texorpdfstring{SHGH Conjecture for Linear Systems on Blown-Up Planes}{SHGH Conjecture for Linear Systems on Blown-Up Planes}}\label{30001233}
\subsection{Definitions and mathematical statement}
For distinct general points $p_1,\ldots,p_r\in{\mathbb{P}}^2_{{\mathbb{C}}}$, let $\mathcal L(d;m_1,\ldots,m_r)$ be the projective space of degree-$d$ curves vanishing to order at least $m_i$ at $p_i$. Its virtual and expected dimensions are
\[
 v=\binom{d+2}{2}-1-\sum_i\binom{m_i+1}{2},
 \qquad e=\max\{-1,v\}.
\]
The system is special if its actual projective dimension exceeds $e$, with the empty system assigned dimension $-1$. On the blowup, a $(-1)$-curve is a smooth rational curve with self-intersection $-1$.

An exact reduced-form version of SHGH is that every Cremona-reduced system is nonspecial: for $m_1\ge m_2\ge\cdots\ge0$ and $d\ge m_1+m_2+m_3$, padding missing multiplicities by zero,
\[
 h^0(\mathcal L)=\max\left\{0,\binom{d+2}{2}
                                  -\sum_i\binom{m_i+1}{2}\right\}.
\]
Here $h^0(\mathcal L)$ means the dimension of the vector space before projectivizing. For each degree/multiplicity choice, general means points in an appropriate nonempty Zariski-open set. The geometric formulation uses $(-1)$-speciality, not an undefined fixed-locus condition for empty systems. The catalogue's URL is a paper by Brambilla--Postinghel, not the authors named in its citation; see \href{https://link.springer.com/article/10.1007/s40574-025-00468-5}{[E1]}.
\par\smallskip\textit{Formulation and concept references:} \href{https://link.springer.com/article/10.1007/s40574-025-00468-5}{[S1]}, \href{https://link.springer.com/article/10.1007/s40574-025-00468-5}{[E1]}.

\subsection{Short English statement}
Are unexpected dimensions of general plane-curve systems entirely explained by exceptional curves, so that every Cremona-reduced system has its expected dimension?
\subsection{Research attempt}
\paragraph{Scope and literature check.}
This attempt by GPT-6 Astra Ultra was prepared on 27 September 2026.
The quantifier is over every finite degree and multiplicity tuple satisfying
the displayed reduced inequality. For each tuple it is enough to produce
one configuration with maximal interpolation rank, since nonvanishing of a
matrix minor is Zariski open. It is not enough to verify finitely many tuples,
or to put every point on a convenient curve and assume that the dimension
has not increased.

The primary text at [E1], Section 3.1, explicitly gives the reduced
nonspeciality formulation and records the conjecture as open. It reports
known cases with at most nine points and with multiplicities at most eleven.
The finite examples checked below have multiplicity at most six and therefore
are verification cases inside established territory, not claimed new cases.
The catalogue's bibliographic correction is retained: the linked article is
by Brambilla--Postinghel. No result in that article is interpreted as a
proof for arbitrary multiplicity and arbitrarily many points.

The approach is to expose the interpolation matrix, construct an elementary
surjectivity certificate in a high-degree range, then test whether restriction
to a line can descend to the sharp reduced range. A critical square matrix
provides a concrete obstruction to a dimension-only version of that descent.
All vector-space dimensions below are before projectivizing. Write
\[
 N_d=\binom{d+2}{2},\quad T=\sum_i\binom{m_i+1}{2},
 \quad A=N_d-T,
\]
with $d\geq0$ and $m_i\geq0$ integral; points of multiplicity zero can be
discarded. The desired value is $h^0=\max(0,A)$.

\paragraph{Lemma 493.1: a precise rank formulation and an exact finite certificate.}
Choose an affine chart containing all the points, with coordinates
$p_i=(a_i,b_i)$. In the basis $x^u y^v$, $u,v\geq0$, $u+v\leq d$,
the jet evaluation matrix has rows $(i,s,t)$ for $s,t\geq0$, $s+t<m_i$,
and entries
\begin{equation}\label{a493:matrix}
 M_{(i,s,t),(u,v)}=
 \begin{cases}
 \binom us\binom vt a_i^{u-s}b_i^{v-t},&u\geq s,\ v\geq t,\\
 0,&\text{otherwise}.
 \end{cases}
\end{equation}
Then $h^0=N_d-\operatorname{rank}M$, and nonspeciality is equivalent to
$\operatorname{rank}M=\min(T,N_d)$. If coordinates are integers and a
$\min(T,N_d)$-square minor is nonzero modulo a prime, that computation
certifies nonspeciality for general complex points for this particular tuple.

\emph{Proof.} Substituting $x=a_i+X$, $y=b_i+Y$ into a polynomial,
the coefficient of $X^sY^t$ is exactly the row in
\eqref{a493:matrix}. Vanishing to order $m_i$ means that all coefficients
with $s+t<m_i$ vanish. There are $\binom{m_i+1}{2}$ such coefficients.
Homogenization identifies polynomials of total degree at most $d$ with
homogeneous plane forms of degree $d$, so the kernel has the asserted
dimension. No factorial division is needed in this matrix.

Every minor is a polynomial with integer coefficients in the point
coordinates. If its value at integer coordinates is nonzero modulo a prime,
the integer value and thus the polynomial itself are nonzero. It remains
nonzero on a nonempty Zariski-open set over $\mathbb C$. Intersecting with
the open condition that points are distinct still gives a nonempty set
when the exhibited points are distinct. Rank cannot exceed
$\min(T,N_d)$ anywhere, so this proves equality on that open set.
Finally the elementary lower bound $h^0\geq\max(0,N_d-T)$ holds for
every configuration and supplies the opposite dimension inequality.
\hfill$\square$

This certificate has a one-sided interpretation that matters in an attempted
disproof. Full rank at one configuration proves generic full rank. A rank
drop at one configuration says nothing by itself about generic rank. To
disprove the conjecture for a tuple one would need a rank deficiency for
general points, equivalently the vanishing of every relevant maximal minor
as a polynomial, not just at chosen coordinates.

\paragraph{Lemma 493.2: explicit interpolation in degree $\sum m_i-1$.}
For arbitrary distinct points, not necessarily general, put
$W=\sum_i m_i>0$. The jet map is surjective whenever $d\geq W-1$.
Thus in this range
\begin{equation}\label{a493:highdegree}
                         h^0=N_d-T.
\end{equation}

\emph{Proof.} Work in an affine chart avoiding all the points at infinity.
For each ordered pair $i\ne j$, choose an affine linear form
$\ell_{ij}$ vanishing at $p_j$ but not at $p_i$; distinctness makes this
possible. Let
\[
                 G_i=\prod_{j\ne i}\ell_{ij}^{m_j},
                 \qquad \deg G_i=W-m_i.
\]
At every $p_j$, $j\ne i$, the polynomial $G_i$ vanishes to order at
least $m_j$. At $p_i$ it has nonzero constant term, so it is invertible in
the finite local algebra
$B_i=\mathbb C[X,Y]/(X,Y)^{m_i}$. For any prescribed jet $q_i\in B_i$,
choose the unique representative $Q_i(X,Y)$ of $q_i/G_i$ whose monomials
have degree less than $m_i$. In original coordinates,
$G_iQ_i(x-a_i,y-b_i)$ has degree at most
$(W-m_i)+(m_i-1)=W-1$. Its jet is $q_i$ at $p_i$ and zero at every
other designated point. Summing over $i$ realizes all prescribed jets
simultaneously. Inclusion in degree $d\geq W-1$ and homogenization prove
surjectivity and \eqref{a493:highdegree}.\hfill$\square$

The bound is sharp for the requirement of \emph{all} distinct configurations.
If all points lie on a line, restriction to that line would have to realize
$W$ independent tangential jet values using only $d+1$ coefficients. Thus
surjectivity is impossible when $d<W-1$. This does not establish sharpness
for general points. For $r$ equal multiplicities $m$, the construction costs
$rm-1$ degrees; reduced form only requires $d\geq3m$ once $r\geq3$.
The size of that gap rules out claiming that this elementary interpolation
argument proves SHGH in the difficult range.

\paragraph{Two completely proved benchmark families.}
First suppose that all nonzero multiplicities are one. For each $d$, take
the triangular grid
\[
                 \Gamma_d=\{(i,j)\in\mathbb Z_{\geq0}^2:i+j\leq d\}.
\]
It has $N_d$ points, and evaluation of degree-at-most-$d$ polynomials on
it is an isomorphism. To prove injectivity, let $F$ vanish on the grid.
Its restriction to $y=0$ has $d+1$ distinct roots and degree at most $d$,
so it is zero and $F=yG$, with $\deg G\leq d-1$. For the remaining
grid points $j\geq1$, $G(i,j)=0$. The translated polynomial
$G(x,y+1)$ vanishes on $\Gamma_{d-1}$ and is zero by induction. The
base case $d=0$ is immediate. Equal dimensions give the isomorphism.
Taking subsets of the grid gives independent rows for $r\leq N_d$;
adding distinct points to the grid gives injectivity for $r\geq N_d$.
Lemma 493.1 therefore proves the desired dimension for every $d$ and every
number of general simple points.

Second, for at most three noncollinear points, a projective coordinate
change sends them to the coordinate vertices. A monomial $X^aY^bZ^c$,
$a+b+c=d$, is allowed exactly when
\[
                 b+c\geq m_1,\qquad a+c\geq m_2,
                 \qquad a+b\geq m_3.
\]
For $m_i\leq d+1$, the monomials violating condition $i$ number
$\binom{m_i+1}{2}$. Two violation sets cannot overlap if
$d\geq m_i+m_j-1$: indeed their simultaneous inequalities would give
$d+\text{(remaining exponent)}\leq m_i+m_j-2$. Consequently under
all three of these degree bounds the number of allowed monomials is exactly
$N_d-T$. Every reduced three-point tuple satisfies these bounds. This also
checks the use of vector dimension, the padding by zero multiplicities,
and the endpoint where the surviving space is zero.

These are elementary derivations of known easy subcases. They suggest that
general position should prevent overlap among conditions. But outside these
families the rows are not simply disjoint coordinate conditions, and that
suggestion does not prove their independence.

\paragraph{The restriction-to-a-line attempt, with its exact loss.}
Fix a configuration with some of its points on a line $L$. Let $I$ be the
indices on the line, let $t=\sum_{i\in I}m_i$, and let $Z'$ be the
residual multiplicity data: decrease $m_i$ by one for $i\in I$ and leave
the others unchanged. For $d\geq1$, restriction gives
\begin{equation}\label{a493:residual}
 h^0(d;Z)\leq h^0(d-1;Z')+\max(0,d+1-t).
\end{equation}

Here is a direct proof, avoiding an assumed surjectivity of restriction.
The restriction of a form to $L\simeq\mathbb P^1$ is a binary form of
degree $d$, vanishing at the distinct selected points to orders at least
$m_i$. Such forms form a space of dimension $\max(0,d+1-t)$, by
factoring the required distinct linear factors. The kernel of restriction
consists of forms divisible by an equation $\ell$ of $L$. Writing
$F=\ell G$, multiplication increases vanishing order by exactly one at
each point on $L$; off $L$, $\ell$ is a unit and the order is unchanged.
Thus that kernel is isomorphic to the residual space of degree $d-1$.
The dimension of the image is at most the dimension of the binary-form
space, proving \eqref{a493:residual}.

Put $A'=N_{d-1}-T'$ for the residual virtual vector dimension and
$b=d+1-t$. Since $T-T'=t$ and $N_d-N_{d-1}=d+1$,
\begin{equation}\label{a493:addvirtual}
                              A=A'+b.
\end{equation}
Even if the residual system is nonspecial, the upper bound in
\eqref{a493:residual} becomes $\max(0,A')+\max(0,b)$. It equals
$\max(0,A)$ when $A'$ and $b$ have the same weak sign. If they have
opposite signs, the excess is exactly
\begin{equation}\label{a493:slack}
                           \min\{|A'|,|b|\}.
\end{equation}
This identity follows by checking which of the two positive quantities is
larger. It identifies the dimension lost by discarding the interaction
between trace and residual conditions.

A good specialized configuration gives an upper bound for the generic
dimension by Lemma 493.1. However, after a specialization has put several
points on one line, the residual points cannot simply be declared general
again. A recursive plan needs either one configuration meeting all the
rank requirements, or a justified deformation argument. An arithmetic
schedule of multiplicity decrements alone does not provide either one.

\paragraph{A sharp failure of the dimension-only induction.}
Consider the reduced tuple
\begin{equation}\label{a493:critical}
                             \mathcal L(19;6^{10}).
\end{equation}
Here $19\geq18$, and $N_{19}=210=T$, so the target is $h^0=0$.
Specialize any $q$ of the ten whole fat points onto a line. Then
$t=6q$, $b=20-6q$, and $A'=6q-20=-b$. If one uses only
\eqref{a493:residual} and an optimal residual dimension, the resulting
upper bound is
\[
                   \max(0,6q-20)+\max(0,20-6q)=|20-6q|.
\]
The minimum over integer $0\leq q\leq10$ is two, attained at $q=3$.
No choice of how many whole points lie on the first line makes this
dimension count prove the required zero. If the residual system is special,
the bound becomes still larger. This disproves the proposed \emph{proof
mechanism}, not SHGH: an exact maximal-rank certificate for
\eqref{a493:critical} is given below.

At $q=3$, the trace has dimension two and the residual virtual dimension
is $-2$. Successful vanishing must use the fact that these apparent trace
sections do not lift to plane curves with all the other prescribed jets.
In the cohomological exact sequence, this is injectivity of a connecting
map from that two-dimensional trace space into residual first cohomology.
The elementary inequality forgets that map. The first unresolved step of
this approach is a general maximal-rank statement for the trace-to-obstruction
map, or a degeneration that computes it. No claim that this map is always
generic is made here.

\paragraph{Finite certificates in the reduced range.}
The companion script works over the prime $65537$ and uses the ten distinct
integer points
\[
\begin{gathered}
 (919,259),(700,968),(997,867),(309,600),(264,364),\\
 (518,940),(632,684),(349,890),(508,45),(681,268).
\end{gathered}
\]
It constructs \eqref{a493:matrix} using binomial coefficients, performs
modular elimination, records the original row and column indices of a
maximal minor, and independently recomputes that minor's determinant.
Columns are ordered by increasing $u$, then increasing $v$; rows are
ordered by the displayed points, then increasing $s$, then increasing
$t$. The rows of each selected minor are in the recorded pivot order.
The resulting certificates are:
\[
\begin{array}{c|r|r|r|r|r}
 (d;m^{10})&N_d&T&\operatorname{rank}M&h^0&\det\bmod65537\\\hline
 (6;2^{10})&28&30&28&0&14321\\
 (7;2^{10})&36&30&30&6&31510\\
 (12;4^{10})&91&100&91&0&48454\\
 (13;4^{10})&105&100&100&5&23435\\
 (15;5^{10})&136&150&136&0&43479\\
 (16;5^{10})&153&150&150&3&19376\\
 (18;6^{10})&190&210&190&0&55207\\
 (19;6^{10})&210&210&210&0&3540\\
 (20;6^{10})&231&210&210&21&31206
\end{array}
\]
Every displayed determinant is nonzero, and the script verifies primality
by trial division up to the square root. In particular the square critical
matrix has determinant $3540$ modulo $65537$. Lemma 493.1 turns each
certificate into a statement for a nonempty Zariski-open set of complex
configurations. No probabilistic rank inference or floating-point threshold
is used. The coordinates were selected with a fixed pseudorandom seed only
to obtain a candidate; the subsequent certificate is exact.

The script and full index lists are retained in
\nolinkurl{_Local/top_problems/continuation_descending_2026_09_27/combinatorics/}
as \texttt{verify.py} and \texttt{results.json}. These computations do not
prove a uniform theorem as $m,r,d$ grow. The size of a finite matrix is not
itself a bound on how large a counterexample tuple would have to be.

\paragraph{Stress tests: speciality, general position, and empty systems.}
Five general points determine a unique smooth conic $Q$. Quartics double
at all five include $Q^2$, while their expected vector dimension is
$15-5\cdot3=0$. In fact the space has dimension one. Restriction of any
such quartic to the conic, parametrized by $\mathbb P^1$, gives a binary
form of degree eight with at least ten zeros counted with multiplicity;
it is identically zero. Thus the quartic is divisible by $Q$. Its residual
conic passes through the five points and is proportional to $Q$. This is
the familiar special system $\mathcal L(4;2^5)$, but $4<2+2+2$, so it
does not contradict the reduced target.

There is also a reduced-looking trap if general position is forgotten.
Take ten distinct points on the smooth projective conic with affine equation
$y=x^2$, for example $(i,i^2)$, $1\leq i\leq10$. For cubics with
simple zeros at these points, $d=3=m_1+m_2+m_3$ and the expected vector
dimension is zero. Nevertheless all products of the conic equation with
a linear form are in the system, giving dimension three. Restriction to
the conic has degree six and ten distinct zeros, so every such cubic is
divisible by the conic; the dimension is exactly three. The script also
checks rank seven for this ten-by-ten cubic matrix. These are special
point configurations, while the triangular-grid argument proves generic
dimension zero for ten simple points. The general-position hypothesis is
therefore essential even in reduced form.

Empty systems are handled throughout by vector dimension zero, corresponding
to projective dimension $-1$. Negative virtual dimension is not a negative
number of actual sections. All modular minors lift to characteristic zero
as integers, so the verification does not import a positive-characteristic
postulation theorem into the complex problem.

\paragraph{Outcome.}
\textbf{UNRESOLVED.} The attempt proves explicit high-degree interpolation
for arbitrary distinct points, proves the simple-point and reduced
three-point benchmark families, and supplies exact finite certificates near
several zero-dimensional thresholds. Its substantive failure analysis is
the unavoidable two-dimensional slack of whole-point line restriction at
$\mathcal L(19;6^{10})$, despite independently certified vanishing there.
The broad reduced range remains unproved and undisproved.

The next concrete target is control of the trace-to-obstruction map rather
than another count of its source and target dimensions. One could test
partial-jet degenerations that split a fat point across a divisor, construct
nonzero minors adapted to those degenerations, and compare their leading
terms at the square thresholds. Each proposed induction would need a
termination argument and a configuration or deformation that keeps every
claimed minor nonzero. None is supplied for arbitrary reduced data.

\subsection{Sources}
\begin{itemize}
\item[S1]Antonio Laface and Luca Ugaglia, Bollettino dell'Unione Matematica Italiana (2025).
\url{https://link.springer.com/article/10.1007/s40574-025-00468-5}.
\textit{Formulation source. Consulted 24 September 2026: Authorship is Brambilla–Postinghel, unlike the imported citation; Conjecture 3.1 is the SHGH statement.}
\item[E1]Maria Chiara Brambilla and Elisa Postinghel, Towards Good Postulation of Fat Points, One Step at a Time, Bollettino dell’Unione Matematica Italiana 18 (2025), 737–749, Sections 1 and 3.1, Conjecture 3.1.
\url{https://link.springer.com/article/10.1007/s40574-025-00468-5}.
\textit{Added primary source: Corrected bibliographic authorship and title for the URL already supplied by the catalogue; dimensions and the Cremona-reduced SHGH formulation. Consulted 24 September 2026.}
\end{itemize}
\end{document}
